\Needspace{21\baselineskip}

\CNsection{1}{Why SDN?}

Traditional networks suffer from fundamental limitations that hinder innovation and agility:

{\def\LTcaptype{none} % do not increment counter
\Needspace{13\baselineskip}
\begin{longtable}[c]{@{}
  >{\raggedright\arraybackslash\hspace{0pt}}m{(\linewidth - 2\tabcolsep) * \real{0.4091}}
  >{\raggedright\arraybackslash\hspace{0pt}}m{(\linewidth - 2\tabcolsep) * \real{0.5909}}@{}}
\toprule\noalign{}
\rowcolor{CNink}\CNth{Problem} & \CNth{Description} \\
\CNheadrulesep
\endhead
\bottomrule\noalign{}
\endlastfoot
\textbf{Distributed Control} & Every device runs its own control plane independently. No single entity has a global view of the network state. \\
\textbf{Vendor Lock-in} & Proprietary hardware + software bundles. Migrating between vendors is costly and complex. \\
\textbf{Slow Provisioning} & Manual, box-by-box configuration via CLI. Weeks to deploy a new service. Human errors cause outages. \\
\textbf{Rigid Architecture} & Adding new protocols requires vendor firmware updates. Innovation cycles measured in years. \\
\textbf{Inconsistent Policies} & Policies must be manually replicated across hundreds of devices. Drift is inevitable. \\
\end{longtable}
}

\begin{CNquote}

\textbf{The Core Insight}: If we separate the ``brain'' (control plane) from the ``muscle'' (data plane), we can centralize intelligence, automate operations, and program the network like software.

\end{CNquote}

\Needspace{14\baselineskip}

\CNsection{2}{Traditional vs SDN Architecture}

\CNchained\CNsubsection{}{Traditional Architecture}

\begin{itemize}
\tightlist
\item
  \textbf{Control plane embedded in every device}
\item
  Each device independently computes forwarding decisions
\item
  Distributed protocols (OSPF, BGP, STP) synchronize state
\item
  Convergence takes time; no global optimization
\end{itemize}

\CNsubsection{}{SDN Architecture}

\begin{itemize}
\tightlist
\item
  \textbf{Control plane extracted and centralized}
\item
  Switches become simple forwarding devices
\item
  Controller has a \textbf{global view} of the entire network
\item
  Programmable via APIs (northbound and southbound)
\end{itemize}

{\def\LTcaptype{none} % do not increment counter
\Needspace{16\baselineskip}
\begin{longtable}[c]{@{}
  >{\raggedright\arraybackslash\hspace{0pt}}m{(\linewidth - 4\tabcolsep) * \real{0.3077}}
  >{\raggedright\arraybackslash\hspace{0pt}}m{(\linewidth - 4\tabcolsep) * \real{0.5000}}
  >{\raggedright\arraybackslash\hspace{0pt}}m{(\linewidth - 4\tabcolsep) * \real{0.1923}}@{}}
\toprule\noalign{}
\rowcolor{CNink}\CNth{Aspect} & \CNth{Traditional} & \CNth{SDN} \\
\CNheadrulesep
\endhead
\bottomrule\noalign{}
\endlastfoot
Control Plane & Distributed (per device) & Centralized (controller) \\
Data Plane & Coupled with control & Separated, simple forwarding \\
Configuration & CLI, box-by-box & API-driven, programmable \\
Vendor & Locked-in & Open, multi-vendor \\
Innovation Speed & Years (firmware cycle) & Days (software development) \\
Visibility & Per-device, fragmented & Global, real-time \\
Automation & Scripts/\allowbreak{}templates & Native API + intent \\
\end{longtable}
}

\CNkeepfig{m15-sdn-layers}{5}

\CNsection{3}{SDN Architecture Layers}

\CNfigure{m15-sdn-layers}{SDN architecture: application, control and infrastructure layers}

\CNsubsection{3.1}{Infrastructure Layer (Data Plane)}

\begin{itemize}
\tightlist
\item
  Physical or virtual \textbf{switches and routers}
\item
  Performs packet forwarding based on \textbf{flow tables} programmed by the controller
\item
  No independent decision-making — follows instructions from the control layer
\item
  Hardware: white-box switches, programmable ASICs, TCAM for fast lookup
\end{itemize}

\CNsubsection{3.2}{Control Layer (SDN Controller)}

\begin{itemize}
\tightlist
\item
  The \textbf{``brain''} of the network
\item
  Maintains a \textbf{global view} of topology, links, hosts, and traffic
\item
  Translates high-level policies into low-level flow rules
\item
  Examples: ONOS, OpenDaylight (ODL), Ryu, Floodlight, Nokia NSP, Cisco ACI
\end{itemize}

\CNsubsection{3.3}{Application Layer}

\begin{itemize}
\tightlist
\item
  \textbf{Business logic and policies} expressed as software applications
\item
  Consumes the controller’s northbound API
\item
  Examples: traffic engineering, security (micro-segmentation), load balancing, monitoring, intent-based networking
\end{itemize}

\Needspace{22\baselineskip}

\CNsection{4}{SDN Controller Functions}

{\def\LTcaptype{none} % do not increment counter
\Needspace{16\baselineskip}
\begin{longtable}[c]{@{}
  >{\raggedright\arraybackslash\hspace{0pt}}m{(\linewidth - 2\tabcolsep) * \real{0.4348}}
  >{\raggedright\arraybackslash\hspace{0pt}}m{(\linewidth - 2\tabcolsep) * \real{0.5652}}@{}}
\toprule\noalign{}
\rowcolor{CNink}\CNth{Function} & \CNth{Description} \\
\CNheadrulesep
\endhead
\bottomrule\noalign{}
\endlastfoot
\textbf{Topology Discovery} & Discovers all network devices, links, and hosts. Builds a graph representation. Uses LLDP, OFDP. \\
\textbf{Path Computation /\allowbreak{} Routing} & Computes optimal paths using global view. Can apply constraints (latency, bandwidth, cost). \\
\textbf{Policy Enforcement} & Translates high-level policies into device-level flow rules. \\
\textbf{API Exposure} & Provides northbound REST/\allowbreak{}gRPC APIs for applications to consume network services. \\
\textbf{State Management} & Maintains consistent network state (flow tables, counters, port status). \\
\textbf{Multi-tenancy} & Network slicing and isolation between tenants. \\
\textbf{Fault Management} & Detects link/\allowbreak{}device failures and reroutes traffic in milliseconds. \\
\end{longtable}
}

\CNsection{5}{Southbound Interfaces}

The southbound interface connects the \textbf{controller to infrastructure devices} (data plane).

\CNsubsection{5.1}{OpenFlow}

\begin{itemize}
\tightlist
\item
  Most well-known SDN southbound protocol
\item
  Defined by \textbf{Open Networking Foundation (ONF)}
\item
  Provides a standard way to program \textbf{flow tables} in switches
\item
  Controller installs/\allowbreak{}modifies/\allowbreak{}deletes flow entries
\item
  Switch reports events (topology changes, packet-in) to controller
\item
  Versions: 1.0 \ensuremath{\to} 1.3 (most deployed) \ensuremath{\to} 1.5
\end{itemize}

\CNsubsection{5.2}{P4 (Programming Protocol-independent Packet Processors)}

\begin{itemize}
\tightlist
\item
  \textbf{Domain-specific language} for programming the data plane itself
\item
  Goes beyond OpenFlow: defines \textbf{what} the switch can match on
\item
  Programmable packet parsing and processing pipelines
\item
  Target-independent: compiles to different hardware (Tofino, FPGA, software switches)
\item
  Use case: custom headers, in-network computing, INT (In-band Network Telemetry)
\end{itemize}

\CNsubsection{5.3}{ForCES (Forwarding and Control Element Separation)}

\begin{itemize}
\tightlist
\item
  IETF standard (RFC 5810)
\item
  Separates Control Element (CE) from Forwarding Element (FE)
\item
  Defines a protocol for CE-FE communication
\item
  Less adopted than OpenFlow in practice
\end{itemize}

\Needspace{15\baselineskip}

\CNsubsection{5.4}{Other Southbound Protocols}

{\def\LTcaptype{none} % do not increment counter
\Needspace{11\baselineskip}
\begin{longtable}[c]{@{}cc@{}}
\toprule\noalign{}
\rowcolor{CNink}\CNth{Protocol} & \CNth{Use Case} \\
\CNheadrulesep
\endhead
\bottomrule\noalign{}
\endlastfoot
NETCONF/\allowbreak{}YANG & Device configuration (model-driven) \\
gNMI/\allowbreak{}gNOI & Streaming telemetry and operations \\
BGP-LS & Topology export from traditional networks to SDN controller \\
PCEP & Path computation for MPLS/\allowbreak{}SR networks \\
\end{longtable}
}

\CNsection{6}{Northbound Interfaces}

The northbound interface connects \textbf{applications to the controller}.

\CNsubsection{6.1}{REST APIs}

\begin{itemize}
\item
  Most common northbound interface
\item
  HTTP-based, JSON/\allowbreak{}XML payloads
\item
  CRUD operations on network resources (flows, topologies, policies)
\item
  Example:

  \tcbset{CNcodemode/.style={unbreakable}}\def\CNcodelang{JSON}

\begin{Shaded}
\begin{Highlighting}[]
\ErrorTok{POST} \ErrorTok{/api/v1/flows}
\FunctionTok{\{}
  \DataTypeTok{"switch"}\FunctionTok{:} \StringTok{"s1"}\FunctionTok{,}
  \DataTypeTok{"match"}\FunctionTok{:} \FunctionTok{\{}\DataTypeTok{"src\_ip"}\FunctionTok{:} \StringTok{"10.0.0.1"}\FunctionTok{,} \DataTypeTok{"dst\_ip"}\FunctionTok{:} \StringTok{"10.0.0.2"}\FunctionTok{\},}
  \DataTypeTok{"action"}\FunctionTok{:} \FunctionTok{\{}\DataTypeTok{"output"}\FunctionTok{:} \StringTok{"port3"}\FunctionTok{\},}
  \DataTypeTok{"priority"}\FunctionTok{:} \DecValTok{100}
\FunctionTok{\}}
\end{Highlighting}
\end{Shaded}
\end{itemize}

\CNsubsection{6.2}{Intent-Based Networking (IBN)}

\begin{itemize}
\tightlist
\item
  Higher abstraction: express \textbf{what} you want, not \textbf{how} to achieve it
\item
  Controller translates intent into device-level configuration
\item
  Example intent: \emph{``Ensure video traffic between Site A and Site B has latency \textless{} 10ms''}
\item
  The controller figures out the path, QoS policies, and failover
\item
  Implementations: ONOS intent framework, Cisco DNA Center, Nokia NSP
\end{itemize}

\CNsubsection{6.3}{gRPC /\allowbreak{} Protobuf}

\begin{itemize}
\tightlist
\item
  High-performance, streaming-capable API
\item
  Used in modern controllers for real-time telemetry and configuration
\item
  Strongly typed (protobuf schemas)
\end{itemize}

\Needspace{14\baselineskip}

\Needspace{25\baselineskip}

\CNsection{7}{OpenFlow Basics}

\CNchained\CNsubsection{7.1}{Flow Table Structure}

Each switch maintains one or more \textbf{flow tables}. Each entry contains:

{\def\LTcaptype{none} % do not increment counter
\Needspace{14\baselineskip}
\begin{longtable}[c]{@{}
  >{\raggedright\arraybackslash\hspace{0pt}}m{(\linewidth - 2\tabcolsep) * \real{0.3500}}
  >{\raggedright\arraybackslash\hspace{0pt}}m{(\linewidth - 2\tabcolsep) * \real{0.6500}}@{}}
\toprule\noalign{}
\rowcolor{CNink}\CNth{Field} & \CNth{Description} \\
\CNheadrulesep
\endhead
\bottomrule\noalign{}
\endlastfoot
\textbf{Match Fields} & Criteria to match packets (src/\allowbreak{}dst MAC, IP, port, VLAN, MPLS label) \\
\textbf{Priority} & Higher priority entries matched first \\
\textbf{Counters} & Packet/\allowbreak{}byte count for statistics \\
\textbf{Instructions/\allowbreak{}Actions} & What to do with matched packets \\
\textbf{Timeouts} & Idle timeout, hard timeout (auto-expire) \\
\textbf{Cookie} & Opaque identifier set by controller \\
\end{longtable}
}

\Needspace{18\baselineskip}

\CNsubsection{7.2}{Actions}

{\def\LTcaptype{none} % do not increment counter
\Needspace{14\baselineskip}
\begin{longtable}[c]{@{}cc@{}}
\toprule\noalign{}
\rowcolor{CNink}\CNth{Action} & \CNth{Description} \\
\CNheadrulesep
\endhead
\bottomrule\noalign{}
\endlastfoot
\CNcode{OUTPUT(port)} & Forward packet to specified port \\
\CNcode{DROP} & Discard the packet \\
\CNcode{SET\_FIELD} & Modify header fields (rewrite MAC, IP, VLAN) \\
\CNcode{PUSH/POP\_VLAN} & Add or remove VLAN tags \\
\CNcode{GROUP} & Send to a group table (multipath, failover) \\
\CNcode{METER} & Rate-limit the flow \\
\end{longtable}
}

\Needspace{17\baselineskip}

\CNsubsection{7.3}{Group Tables}

Enable advanced forwarding:

{\def\LTcaptype{none} % do not increment counter
\Needspace{11\baselineskip}
\begin{longtable}[c]{@{}cc@{}}
\toprule\noalign{}
\rowcolor{CNink}\CNth{Group Type} & \CNth{Use Case} \\
\CNheadrulesep
\endhead
\bottomrule\noalign{}
\endlastfoot
\textbf{All} & Multicast/\allowbreak{}broadcast (copy to all buckets) \\
\textbf{Select} & Load balancing (round-robin or weighted) \\
\textbf{Indirect} & Single next-hop (efficient pointer) \\
\textbf{Fast Failover} & First live bucket used (link protection) \\
\end{longtable}
}

\CNsubsection{7.4}{Meter Tables}

\begin{itemize}
\tightlist
\item
  Rate limiting per flow
\item
  Defined in \textbf{bands}: rate threshold + action (drop or DSCP remark)
\item
  Enables QoS enforcement at the switch level
\end{itemize}

\CNkeepfig{m15-pipeline}{3}

\CNsubsection{7.5}{Packet Processing Pipeline}

\CNfigure{m15-pipeline}{OpenFlow packet-processing pipeline}

\Needspace{14\baselineskip}

\CNsection{8}{SDN in Telecom}

\CNchained\CNsubsection{8.1}{Transport SDN (WAN Optimization)}

\begin{itemize}
\tightlist
\item
  Centralized control of optical/\allowbreak{}MPLS transport networks
\item
  Dynamic bandwidth allocation based on traffic demand
\item
  Multi-layer optimization (IP over optical)
\item
  Examples: Nokia NSP, Ciena Blue Planet, Cisco WAE
\end{itemize}

\CNsubsection{8.2}{SD-WAN (Software-Defined Wide Area Network)}

\begin{itemize}
\tightlist
\item
  SDN principles applied to enterprise WAN
\item
  Centralized controller manages branch CPE devices
\item
  Overlays on top of multiple underlay transports (MPLS, Internet, LTE)
\item
  Application-aware routing (route voice over MPLS, web over Internet)
\item
  Zero-touch provisioning for branch sites
\item
  Vendors: Cisco Viptela, VMware VeloCloud, Fortinet, Versa
\end{itemize}

\CNkeepfig{m15-smf-upf}{3}

\CNsubsection{8.3}{5G Integration: SMF as SDN Controller for UPF}

\CNfigure{m15-smf-upf}{The SMF acting as an SDN controller for the UPFs}

\Needspace{16\baselineskip}

\textbf{The 5G User Plane is inherently SDN-like:}

{\def\LTcaptype{none} % do not increment counter
\Needspace{13\baselineskip}
\begin{longtable}[c]{@{}
  >{\raggedright\arraybackslash\hspace{0pt}}m{(\linewidth - 2\tabcolsep) * \real{0.4643}}
  >{\raggedright\arraybackslash\hspace{0pt}}m{(\linewidth - 2\tabcolsep) * \real{0.5357}}@{}}
\toprule\noalign{}
\rowcolor{CNink}\CNth{SDN Concept} & \CNth{5G Equivalent} \\
\CNheadrulesep
\endhead
\bottomrule\noalign{}
\endlastfoot
SDN Controller & \textbf{SMF} (Session Management Function) \\
Data Plane Switch & \textbf{UPF} (User Plane Function) \\
Southbound Protocol & \textbf{PFCP} (Packet Forwarding Control Protocol) \\
Flow Rules & \textbf{PDRs} (Packet Detection Rules) + \textbf{FARs} (Forwarding Action Rules) \\
QoS Enforcement & \textbf{QERs} (QoS Enforcement Rules) \\
\end{longtable}
}

\CNsubsection{8.4}{UPF Programmability via PFCP}

PFCP (3GPP TS 29.244) is the protocol between SMF and UPF:

\begin{itemize}
\tightlist
\item
  \textbf{PDR (Packet Detection Rule)}: Match criteria (source interface, IP, GTP TEID)
\item
  \textbf{FAR (Forwarding Action Rule)}: Action (forward, drop, buffer, encapsulate)
\item
  \textbf{QER (QoS Enforcement Rule)}: Rate limiting (MBR, GBR)
\item
  \textbf{URR (Usage Reporting Rule)}: Traffic measurement and reporting
\end{itemize}

This is functionally equivalent to OpenFlow’s match-action model!

\Needspace{22\baselineskip}

\CNsection{9}{Benefits of SDN}

{\def\LTcaptype{none} % do not increment counter
\Needspace{16\baselineskip}
\begin{longtable}[c]{@{}
  >{\raggedright\arraybackslash\hspace{0pt}}m{(\linewidth - 2\tabcolsep) * \real{0.4091}}
  >{\raggedright\arraybackslash\hspace{0pt}}m{(\linewidth - 2\tabcolsep) * \real{0.5909}}@{}}
\toprule\noalign{}
\rowcolor{CNink}\CNth{Benefit} & \CNth{Description} \\
\CNheadrulesep
\endhead
\bottomrule\noalign{}
\endlastfoot
\textbf{Programmability} & Network is software-controlled; automate everything via APIs \\
\textbf{Automation} & Zero-touch provisioning, auto-remediation, CI/\allowbreak{}CD for network \\
\textbf{Multi-vendor} & Open interfaces allow mixing hardware vendors (white-box revolution) \\
\textbf{Fast Innovation} & New services deployed in hours/\allowbreak{}days vs. months \\
\textbf{Global Optimization} & Centralized view enables optimal path computation \\
\textbf{Network Slicing} & Dynamic resource allocation per tenant/\allowbreak{}service \\
\textbf{Cost Reduction} & Commodity hardware + open-source software vs. proprietary boxes \\
\end{longtable}
}

\Needspace{14\baselineskip}

\CNsection{10}{Challenges of SDN}

\CNchained\CNsubsection{10.1}{Controller Scalability}

\begin{itemize}
\tightlist
\item
  Single controller cannot handle thousands of switches + millions of flows
\item
  \textbf{Solution}: Controller clustering (ONOS, ODL support distributed controllers)
\item
  Trade-off: consistency vs. availability (CAP theorem applies)
\end{itemize}

\CNsubsection{10.2}{Single Point of Failure}

\begin{itemize}
\tightlist
\item
  If the controller fails, no new flows can be installed
\item
  Data plane continues forwarding existing flows (proactive mode)
\item
  \textbf{Mitigation}: Controller redundancy (active-standby, active-active)
\end{itemize}

\CNsubsection{10.3}{East-West Interface}

\begin{itemize}
\tightlist
\item
  Communication between controller instances in a cluster
\item
  No standard protocol (each controller has proprietary clustering)
\item
  Challenges: state synchronization, split-brain scenarios, consistency
\end{itemize}

\Needspace{15\baselineskip}

\CNsubsection{10.4}{Security}

{\def\LTcaptype{none} % do not increment counter
\Needspace{11\baselineskip}
\begin{longtable}[c]{@{}
  >{\raggedright\arraybackslash\hspace{0pt}}m{(\linewidth - 2\tabcolsep) * \real{0.5000}}
  >{\raggedright\arraybackslash\hspace{0pt}}m{(\linewidth - 2\tabcolsep) * \real{0.5000}}@{}}
\toprule\noalign{}
\rowcolor{CNink}\CNth{Threat} & \CNth{Impact} \\
\CNheadrulesep
\endhead
\bottomrule\noalign{}
\endlastfoot
Controller compromise & Attacker controls entire network \\
Man-in-the-middle (controller\ensuremath{\leftrightarrow}switch) & Flow rule injection/\allowbreak{}modification \\
DoS on controller & Overwhelm with Packet-In messages \\
Malicious application & Rogue app installs harmful flows via NBI \\
\end{longtable}
}

\textbf{Mitigations}: TLS on OpenFlow channel, RBAC for northbound API, rate-limiting Packet-In, application sandboxing.

\CNboxsection{Key Takeaways}

\begin{CNbox}[unbreakable]{sum}{Key Takeaways}

\begin{enumerate}
\def\labelenumi{\arabic{enumi}.}
\tightlist
\item
  SDN \textbf{separates} the control plane from the data plane
\item
  A \textbf{centralized controller} provides a global view and programmability
\item
  \textbf{OpenFlow} is the foundational southbound protocol (match-action paradigm)
\item
  \textbf{P4} extends programmability to the data plane pipeline itself
\item
  5G architecture is \textbf{inherently SDN-like} (SMF controls UPF via PFCP)
\item
  SDN enables \textbf{network automation, slicing, and rapid service deployment}
\item
  Challenges remain around \textbf{scalability, security, and standardization}
\end{enumerate}

\end{CNbox}
